\section{Stable enclosures and certification of a returned policy}\label{sec:stability61}
A valid one-step interval and a stable recursion are different requirements. The first establishes a runtime bracket; the second controls how an approximation allowance propagates when the next-date input changes. We make the distinction explicit and then apply the common-reference construction to a policy that has already been returned by a prospective service.

\subsection{Fixed-cover stability}
At a fixed approximation level, hold the state cells, action covers, candidate actions, innovation bins and their probabilities fixed. An endpoint map may take minima or maxima over the cells intersecting a transition box, average those endpoints with nonnegative weights of total mass one, add a stage-cost endpoint and finally minimize over a fixed finite family of action intervals or feasible candidates. A reference-shifted endpoint map may additionally add a fixed signed reference advantage. These operations describe the exact endpoint maps beneath the outward arithmetic used here.

\begin{lemma}[Stability of the executable endpoint maps]\label{lem:stable61}
Every such exact endpoint map $\mathcal E_t$ is order preserving and satisfies
\begin{equation}\label{eq:endpoint-stability61}
 \|\mathcal E_t f-\mathcal E_t g\|_\infty
       \le\beta_t\|f-g\|_\infty.
\end{equation}
Before intersection with an independent bound, it also satisfies
$\mathcal E_t(f+c)=\mathcal E_t f+\beta_t c$ for constants $c$. Intersecting a lower endpoint with an independently valid lower bound by taking a maximum, or an upper endpoint with an independently valid upper bound by taking a minimum, preserves order and the nonexpansive estimate when the independent bound is fixed. A numerical endpoint perturbation with verified uniform magnitude $\delta_t$ enters the propagated error as an additional local allowance, rather than being presumed nonexpansive because it is outward.
\end{lemma}
\begin{proof}
For any nonempty finite index set $I$, both $\min_{i\in I}f_i$ and $\max_{i\in I}f_i$ are order preserving and one-Lipschitz in the supremum norm. A nonnegative weighted average of total mass one has the same properties. A fixed additive term does not alter a difference, and multiplication by $\beta_t\ge0$ supplies the factor in \eqref{eq:endpoint-stability61}. The final minimum over actions is again one-Lipschitz. Each of these operations commutes with addition of a constant in the stated way. Finally, $x\mapsto\min\{x,c\}$ and $x\mapsto\max\{x,c\}$ are monotone and one-Lipschitz. A separate numerical perturbation is bounded by the triangle inequality.
\end{proof}

Consequently, the consistency clause of Theorem~\ref{thm:bellman-bracket60} applies to the fixed-cover maps actually used in the computation. If the ideal lower and upper maps at the true next Bellman value have errors at most $\alpha_t$, and their numerical evaluation has uniformly bounded endpoint error $\delta_t$ on the relevant range of inputs, put $\eta_t=\alpha_t+\delta_t$. Then the two one-sided errors satisfy
\[
 e_t^-\le\eta_t+\beta_te_{t+1}^- ,\qquad
 e_t^+\le\eta_t+\beta_te_{t+1}^+ .
\]
Their sum gives the displayed two-sided finite-horizon bound. Different approximation levels may use different covers; the stability comparison is within each fixed level. Arbitrarily small tolerances require all nonvanishing sources of error to be refined, including the action quantum when it is a binding discretization. A finite ladder at quantum $1/4096$ is not, by itself, an executed zero-error limit.

Inclusion alone would not justify this step. In a one-state example with continuation $v\ge0$ and discount $\beta$, the endpoint $\beta v+v^2$ is a valid upper bound on $\beta v$ and is exact at $v=0$. Its variation from $0$ to $1$ nevertheless exceeds $\beta$. Thus an ideal error statement only at the true continuation cannot replace \eqref{eq:endpoint-stability61}. The present proof supplies the needed property for the actual min--max and probability-average maps; it does not assume stability of every conceivable enclosing procedure.

\subsection{The original optimum and the actual returned actor}
Let $h_t=J_t^0$ denote the value of the installed zero-investment policy in the original model. A common-reference Bellman calculation supplies a function $\ell_t$ satisfying
\[
 \ell_t(x)\le J_t^*(x)-h_t(x).
\]
Now fix a complete, previously returned acquired policy $\pi$. Its observation partition and quantum action indices are part of its identity and are not changed by verification. On a verification cell $C$, let $R_t(C)$ contain every action actually taken by $\pi_t$ at a state in $C$. Such a rectangle is an outer enclosure of the graph of the fixed policy, not a proposal to implement every action in that rectangle. In particular, its endpoints need not be feasible throughout $C$; the actual policy must remain feasible at each actual state, and the primitive formulas used on the outer enclosure must remain valid there.

Suppose $\mathcal V_t(C)$ is an outward enclosure, for all actual graph points $(x,\pi_t(x))$ with $x\in C$, of
\begin{equation}\label{eq:fixed-excess61}
 A_t^0(x,\pi_t(x))+
       \beta_t P_t^{\pi_t(x)}u_{t+1}(x),
\end{equation}
where $A_t^0(x,a)=c_t(x,a)+\beta_tP_t^a h_{t+1}(x)-h_t(x)$ and $u_T=0$. Define $\widetilde u_t(C)$ as the upper endpoint of $\mathcal V_t(C)$.

\begin{proposition}[Paid revalidation of an unchanged policy]\label{prop:fixed-policy61}
If the stored actor is feasible and the enclosures in \eqref{eq:fixed-excess61} include every reached state and the original innovation law, the recursion $u_t=\widetilde u_t$ gives
\[
 \ell_t(x)\le J_t^*(x)-h_t(x)
       \le J_t^\pi(x)-h_t(x)\le u_t(x).
\]
If an independent whole-domain policy certificate has established $J_t^\pi\le h_t$ at every date, the recursion may instead use
\begin{equation}\label{eq:monotone-intersection61}
 u_t(C)=\min\{\widetilde u_t(C),0\}.
\end{equation}
In either case, an outward maximum of $u_t(C)-\ell_t(C)$ bounds the loss of this particular returned actor against the original continuous-action Bellman optimum. Verification does not change the actor or require a policy-cost sampling event.
\end{proposition}
\begin{proof}
The identity
\[
 J_t^\pi-h_t=A_t^0(\cdot,\pi_t)
       +\beta_tP_t^{\pi_t}(J_{t+1}^\pi-h_{t+1})
\]
holds under the original kernels. Starting from the common terminal value, monotonicity of expectation and the whole-graph enclosure prove the upper inequality by backward induction. The lower inequality is the previously established Bellman lower bound; feasibility gives the middle inequality. When $J_t^\pi-h_t\le0$ has been independently proved, intersection with zero preserves the upper bound at each induction step. Subtracting the two bounds cancels the same $h_t(x)$ pointwise. No difference between extrema at unrelated states is substituted for this cancellation.
\end{proof}

The independent property used in \eqref{eq:monotone-intersection61} is obtained here by reconstructing the original reference-advantage gate at every acquired cell and date. Its nonpositive upper endpoints imply $A_t^0(x,\pi_t(x))\le0$ throughout the domain, and backward policy improvement gives $J_t^\pi\le J_t^0$. The revalidation also retains the entirely unclipped upper recursion and its final gaps. This permits the reader to distinguish a justified independent bound from a numerically favorable truncation. A record of nonpositive values at sampled centers would not supply the required property.

The verification grid is finer than the actor's acquisition partition. Actor ranges include all leaves intersecting each closed verification cell, including both sides of a discontinuity. Rectangle extrema and original-law innovation bins then enclose the continuation. These calculations have a real cost: the common lower-table construction, each fixed-actor upper recursion and the whole-process return overhead are separate measured services. They do not retroactively reduce a prospective service's original stopping cost or certify an unexecuted high-dimensional case.
